\documentclass[12pt]{article}
\usepackage[utf8]{inputenc}
\usepackage[spanish]{babel}
\usepackage{geometry}
\usepackage{amsmath}

\geometry{a4paper, margin=1in}

\title{Plan de Estudio de Física: Semana 2 \\ \large Selección de Problemas}
\author{Basado en 'Física, 6ta Edición' de Wilson, Buffa y Lou}
\date{\today}

\begin{document}

\maketitle

\section*{Problemas Recomendados de la Semana 2}

Esta selección de 10 problemas, extraídos del libro de texto, cubre los temas de 'Movimiento Avanzado y Leyes de Newton' (Capítulos 3 y 4). Los ejercicios están ordenados por dificultad creciente para facilitar un aprendizaje progresivo.

\subsection*{Nivel Básico (Conceptos Fundamentales)}

\begin{enumerate}
    \item \textbf{Capítulo 3, Ejercicio 22 (página 96):}
    \begin{itemize}
        \item[\textbf{Tema:}] Suma de vectores.
        \item[\textbf{Descripción:}] Este problema es ideal para empezar. Te pide encontrar la resultante de la suma de dos vectores, lo que te permitirá practicar los conceptos básicos de adición de vectores de manera gráfica y analítica.
    \end{itemize}

    \item \textbf{Capítulo 4, Ejercicio 1 (página 131):}
    \begin{itemize}
        \item[\textbf{Tema:}] Conceptos de fuerza y fuerza neta.
        \item[\textbf{Descripción:}] Es una pregunta conceptual que te ayudará a diferenciar entre masa y peso, y a entender cómo la fuerza neta afecta el estado de movimiento de un objeto.
    \end{itemize}

    \item \textbf{Capítulo 4, Ejercicio 15 (página 132):}
    \begin{itemize}
        \item[\textbf{Tema:}] Segunda Ley de Newton.
        \item[\textbf{Descripción:}] Un excelente ejercicio para aplicar directamente la Segunda Ley de Newton ($F=ma$). Involucra dos fuerzas actuando sobre un objeto, requiriendo que calcules la fuerza neta para encontrar la aceleración.
    \end{itemize}
\end{enumerate}

\begin{enumerate}
\setcounter{enumi}{3}
    \item \textbf{Capítulo 3, Ejercicio 55 (página 98):}
    \begin{itemize}
        \item[\textbf{Tema:}] Movimiento de proyectiles (lanzamiento horizontal).
        \item[\textbf{Descripción:}] Este es un problema clásico de un proyectil lanzado horizontalmente. Te permitirá practicar cómo descomponer el movimiento en sus componentes horizontal (velocidad constante) y vertical (caída libre).
    \end{itemize} 

    \item \textbf{Capítulo 4, Ejercicio 30 (página 133):}
    \begin{itemize}
        \item[\textbf{Tema:}] Segunda Ley de Newton y Fricción.
        \item[\textbf{Descripción:}] Introduce la fuerza de fricción. Deberás aplicar la Segunda Ley de Newton para determinar la aceleración de un objeto cuando se le aplica una fuerza neta y hay fricción oponiéndose al movimiento.
    \end{itemize} 

    \item \textbf{Capítulo 3, Ejercicio 86 (página 101):}
    \begin{itemize}
        \item[\textbf{Tema:}] Velocidad relativa.
        \item[\textbf{Descripción:}] Un problema introductorio sobre velocidad relativa en una dimensión. Te ayudará a comprender cómo sumar o restar velocidades dependiendo de los marcos de referencia.
    \end{itemize} 

    \item \textbf{Capítulo 4, Ejercicio 57 (página 135):}
    \begin{itemize}
        \item[\textbf{Tema:}] Diagramas de cuerpo libre y equilibrio.
        \item[\textbf{Descripción:}] Este problema requiere que dibujes un diagrama de cuerpo libre para un objeto en un plano inclinado y que apliques las condiciones de equilibrio. Es fundamental para dominar el análisis de fuerzas.
    \end{itemize} 
\end{enumerate}

\begin{enumerate}
\setcounter{enumi}{7}
    \item \textbf{Capítulo 3, Ejercicio 75 (página 100):}
    \begin{itemize}
        \item[\textbf{Tema:}] Movimiento de proyectiles (lanzamiento con ángulo).
        \item[\textbf{Descripción:}] Un problema más completo de movimiento de proyectiles, donde el lanzamiento ocurre desde una altura y con un ángulo. Aquí deberás usar las ecuaciones cinemáticas para ambas componentes (x e y) para encontrar el alcance.
    \end{itemize}

    \item \textbf{Capítulo 4, Ejercicio 71 (página 135):}
    \begin{itemize}
        \item[\textbf{Tema:}] Sistema de dos cuerpos y plano inclinado.
        \item[\textbf{Descripción:}] Este es un problema complejo que involucra un sistema de dos masas conectadas por una cuerda, con una de ellas en un plano inclinado con fricción. Deberás aplicar la Segunda Ley de Newton a cada masa por separado para resolverlo.
    \end{itemize}

    \item \textbf{Capítulo 3, Ejercicio 101 (página 101):}
    \begin{itemize}
        \item[\textbf{Tema:}] Velocidad relativa en dos dimensiones.
        \item[\textbf{Descripción:}] Es un problema desafiante sobre velocidad relativa en 2D (un nadador cruzando un río). Requiere una cuidadosa suma de vectores para encontrar la velocidad resultante y determinar la trayectoria y el tiempo para cruzar.
    \end{itemize}
\end{enumerate}

\end{document}